% =====================================================================
%  Chapter 7 — Conclusion
%  Input with:  % =====================================================================
%  Chapter 7 — Conclusion
%  Input with:  % =====================================================================
%  Chapter 7 — Conclusion
%  Input with:  % =====================================================================
%  Chapter 7 — Conclusion
%  Input with:  \input{chapters/07_conclusion}
%  No tables or figures. No new citations beyond those already used.
% =====================================================================

\section{Conclusion}
\label{sec:conclusion}

\subsection{The two answers}
\label{sec:two-answers}

Fertility in the post-Soviet space did not converge between 2000 and 2023. The four Central Asian countries remained above the other ten in every year of the period without a single crossing, initially higher-fertility countries did not decline faster, and the difference between the two groups ended wider than it began. Beneath a pooled dispersion series that appears flat, Central Asia grew more homogeneous while the rest of the region pulled apart. Around 2016 the pattern changed: dispersion reached its minimum and reversed, and the difference between the groups moved from 1.19 children per woman to 1.65 in seven years. Roughly seventy per cent of that widening came from the comparison countries falling rather than from Central Asia climbing.

The second answer is bounded, in the way set out in Chapter~\ref{sec:introduction}. The Central Asian difference is $+1.314$ children per woman before controls and $+0.964$ after them. Income, urbanisation, remittances and child survival account for about a quarter of it. The remaining three quarters is not explained by this analysis, and the analysis is clear about why: within countries these variables show no stable association with fertility, and across countries the markers that do correlate with the difference --- Muslim population share, age at marriage, female schooling --- cannot be separated from Central Asian regional identity in a sample of fourteen. What the thesis establishes is the size and durability of what remains, not its composition.

\subsection{What is new here}
\label{sec:contribution}

Three things.

The first is coverage. No previous study places all fourteen comparable post-Soviet states in a single frame with consistent measurement, a common window and inference appropriate to the sample size. The comparative post-communist literature largely excludes Central Asia from its formal analysis; the Central Asian literature compares within the region. This thesis puts both halves of the former Soviet Union side by side, and the fact that the difference never closes in any year of the panel is only visible from that vantage point.

The second is timing. The reversal around 2016 and the widening that follows it postdate the windows of essentially all the comparative work on the region. So does the Uzbek fertility surge, the largest single movement in the panel, which took that country from 2.57 to 3.50 children per woman in six years while its population of women in peak childbearing ages was shrinking. The finding that the recent divergence comes mainly from the low-fertility side rather than the high-fertility side is, as far as I am aware, not documented elsewhere, and it changes what needs explaining.

The third is the boundary itself. A null result reported carefully is worth more than a positive result reported loosely, and this thesis reports several. No control variable is individually distinguishable from zero. Within-country variation yields nothing stable. The cross-section cannot separate religion from region. Rather than working around these, the analysis treats them as the finding: the standard macroeconomic account of post-communist fertility, applied with four covariates over twenty-four years, does not reach the Central Asian difference. That is informative precisely because the account is the one the comparative literature relies on.

Some of the methodological choices were made in the same spirit. The wild-cluster bootstrap returns a p-value of 0.033 on the primary estimate where the asymptotic clustered p-value is below 0.01, and the thesis reports the former. Eight missing remittance observations were left missing rather than interpolated, seven of them in the group under study. And an earlier version of the projection-sensitivity check produced a convergence result that turned out to be an artefact of a shrinking country sample; rebuilding it on balanced samples reversed the finding.

\subsection{What follows}
\label{sec:further-research}

The clearest implication is about units of analysis. This thesis shows that the Central Asian difference lives in persistent between-country variation and that a country-level design cannot decompose it. The literature that can is micro-level, and it is fragmented: separate surveys, separate countries, separate periods. A harmonised individual-level dataset covering the Central Asian states and a comparison group elsewhere in the former Soviet Union, with ethnicity, parity, marital history and religiosity measured consistently, would allow the mechanisms discussed in Chapter~\ref{sec:discussion} to be tested against one another rather than assembled from separate studies.

Three narrower directions follow from specific limitations. Because period fertility rates confound the number of births with their timing, tempo-adjusted or cohort measures would establish how much of the recent divergence is a change in family size and how much a change in scheduling --- a question this thesis raises repeatedly and cannot settle. Because ethnic composition may account for the majority of the Central Asian fertility increase, and this design has no ethnic dimension at all, that omission is the most consequential one to repair. And the generational reading offered in Chapter~\ref{sec:discussion} --- that the Soviet-era dampening of the pathway from religiosity to fertility weakens for cohorts formed after 1991, and that these are the cohorts now at peak childbearing --- is a conjecture that can be tested directly by interacting religiosity with birth cohort in existing survey data. It is the only mechanism I found that predicts the timing of the divergence rather than just its level, which makes it worth testing.

\subsection{Closing}
\label{sec:closing}

The countries in this study spent most of the twentieth century under one government, and much of the literature on their demography has been organised around the expectation that they would grow more alike once it ended. On fertility they have not. Two of them entered this century with rates 1.4 children apart and left it 2.3 apart, and the widening is recent enough that most of the research on the region predates it.

This thesis measures that divergence, shows that the standard economic explanations do not account for the greater part of it, and marks out the boundary of what fourteen countries observed over twenty-four years can establish. What lies beyond the boundary --- marriage, ethnicity, gender norms, expectations about the future --- is where the answer is, and it is not reachable from here.

%  No tables or figures. No new citations beyond those already used.
% =====================================================================

\section{Conclusion}
\label{sec:conclusion}

\subsection{The two answers}
\label{sec:two-answers}

Fertility in the post-Soviet space did not converge between 2000 and 2023. The four Central Asian countries remained above the other ten in every year of the period without a single crossing, initially higher-fertility countries did not decline faster, and the difference between the two groups ended wider than it began. Beneath a pooled dispersion series that appears flat, Central Asia grew more homogeneous while the rest of the region pulled apart. Around 2016 the pattern changed: dispersion reached its minimum and reversed, and the difference between the groups moved from 1.19 children per woman to 1.65 in seven years. Roughly seventy per cent of that widening came from the comparison countries falling rather than from Central Asia climbing.

The second answer is bounded, in the way set out in Chapter~\ref{sec:introduction}. The Central Asian difference is $+1.314$ children per woman before controls and $+0.964$ after them. Income, urbanisation, remittances and child survival account for about a quarter of it. The remaining three quarters is not explained by this analysis, and the analysis is clear about why: within countries these variables show no stable association with fertility, and across countries the markers that do correlate with the difference --- Muslim population share, age at marriage, female schooling --- cannot be separated from Central Asian regional identity in a sample of fourteen. What the thesis establishes is the size and durability of what remains, not its composition.

\subsection{What is new here}
\label{sec:contribution}

Three things.

The first is coverage. No previous study places all fourteen comparable post-Soviet states in a single frame with consistent measurement, a common window and inference appropriate to the sample size. The comparative post-communist literature largely excludes Central Asia from its formal analysis; the Central Asian literature compares within the region. This thesis puts both halves of the former Soviet Union side by side, and the fact that the difference never closes in any year of the panel is only visible from that vantage point.

The second is timing. The reversal around 2016 and the widening that follows it postdate the windows of essentially all the comparative work on the region. So does the Uzbek fertility surge, the largest single movement in the panel, which took that country from 2.57 to 3.50 children per woman in six years while its population of women in peak childbearing ages was shrinking. The finding that the recent divergence comes mainly from the low-fertility side rather than the high-fertility side is, as far as I am aware, not documented elsewhere, and it changes what needs explaining.

The third is the boundary itself. A null result reported carefully is worth more than a positive result reported loosely, and this thesis reports several. No control variable is individually distinguishable from zero. Within-country variation yields nothing stable. The cross-section cannot separate religion from region. Rather than working around these, the analysis treats them as the finding: the standard macroeconomic account of post-communist fertility, applied with four covariates over twenty-four years, does not reach the Central Asian difference. That is informative precisely because the account is the one the comparative literature relies on.

Some of the methodological choices were made in the same spirit. The wild-cluster bootstrap returns a p-value of 0.033 on the primary estimate where the asymptotic clustered p-value is below 0.01, and the thesis reports the former. Eight missing remittance observations were left missing rather than interpolated, seven of them in the group under study. And an earlier version of the projection-sensitivity check produced a convergence result that turned out to be an artefact of a shrinking country sample; rebuilding it on balanced samples reversed the finding.

\subsection{What follows}
\label{sec:further-research}

The clearest implication is about units of analysis. This thesis shows that the Central Asian difference lives in persistent between-country variation and that a country-level design cannot decompose it. The literature that can is micro-level, and it is fragmented: separate surveys, separate countries, separate periods. A harmonised individual-level dataset covering the Central Asian states and a comparison group elsewhere in the former Soviet Union, with ethnicity, parity, marital history and religiosity measured consistently, would allow the mechanisms discussed in Chapter~\ref{sec:discussion} to be tested against one another rather than assembled from separate studies.

Three narrower directions follow from specific limitations. Because period fertility rates confound the number of births with their timing, tempo-adjusted or cohort measures would establish how much of the recent divergence is a change in family size and how much a change in scheduling --- a question this thesis raises repeatedly and cannot settle. Because ethnic composition may account for the majority of the Central Asian fertility increase, and this design has no ethnic dimension at all, that omission is the most consequential one to repair. And the generational reading offered in Chapter~\ref{sec:discussion} --- that the Soviet-era dampening of the pathway from religiosity to fertility weakens for cohorts formed after 1991, and that these are the cohorts now at peak childbearing --- is a conjecture that can be tested directly by interacting religiosity with birth cohort in existing survey data. It is the only mechanism I found that predicts the timing of the divergence rather than just its level, which makes it worth testing.

\subsection{Closing}
\label{sec:closing}

The countries in this study spent most of the twentieth century under one government, and much of the literature on their demography has been organised around the expectation that they would grow more alike once it ended. On fertility they have not. Two of them entered this century with rates 1.4 children apart and left it 2.3 apart, and the widening is recent enough that most of the research on the region predates it.

This thesis measures that divergence, shows that the standard economic explanations do not account for the greater part of it, and marks out the boundary of what fourteen countries observed over twenty-four years can establish. What lies beyond the boundary --- marriage, ethnicity, gender norms, expectations about the future --- is where the answer is, and it is not reachable from here.

%  No tables or figures. No new citations beyond those already used.
% =====================================================================

\section{Conclusion}
\label{sec:conclusion}

\subsection{The two answers}
\label{sec:two-answers}

Fertility in the post-Soviet space did not converge between 2000 and 2023. The four Central Asian countries remained above the other ten in every year of the period without a single crossing, initially higher-fertility countries did not decline faster, and the difference between the two groups ended wider than it began. Beneath a pooled dispersion series that appears flat, Central Asia grew more homogeneous while the rest of the region pulled apart. Around 2016 the pattern changed: dispersion reached its minimum and reversed, and the difference between the groups moved from 1.19 children per woman to 1.65 in seven years. Roughly seventy per cent of that widening came from the comparison countries falling rather than from Central Asia climbing.

The second answer is bounded, in the way set out in Chapter~\ref{sec:introduction}. The Central Asian difference is $+1.314$ children per woman before controls and $+0.964$ after them. Income, urbanisation, remittances and child survival account for about a quarter of it. The remaining three quarters is not explained by this analysis, and the analysis is clear about why: within countries these variables show no stable association with fertility, and across countries the markers that do correlate with the difference --- Muslim population share, age at marriage, female schooling --- cannot be separated from Central Asian regional identity in a sample of fourteen. What the thesis establishes is the size and durability of what remains, not its composition.

\subsection{What is new here}
\label{sec:contribution}

Three things.

The first is coverage. No previous study places all fourteen comparable post-Soviet states in a single frame with consistent measurement, a common window and inference appropriate to the sample size. The comparative post-communist literature largely excludes Central Asia from its formal analysis; the Central Asian literature compares within the region. This thesis puts both halves of the former Soviet Union side by side, and the fact that the difference never closes in any year of the panel is only visible from that vantage point.

The second is timing. The reversal around 2016 and the widening that follows it postdate the windows of essentially all the comparative work on the region. So does the Uzbek fertility surge, the largest single movement in the panel, which took that country from 2.57 to 3.50 children per woman in six years while its population of women in peak childbearing ages was shrinking. The finding that the recent divergence comes mainly from the low-fertility side rather than the high-fertility side is, as far as I am aware, not documented elsewhere, and it changes what needs explaining.

The third is the boundary itself. A null result reported carefully is worth more than a positive result reported loosely, and this thesis reports several. No control variable is individually distinguishable from zero. Within-country variation yields nothing stable. The cross-section cannot separate religion from region. Rather than working around these, the analysis treats them as the finding: the standard macroeconomic account of post-communist fertility, applied with four covariates over twenty-four years, does not reach the Central Asian difference. That is informative precisely because the account is the one the comparative literature relies on.

Some of the methodological choices were made in the same spirit. The wild-cluster bootstrap returns a p-value of 0.033 on the primary estimate where the asymptotic clustered p-value is below 0.01, and the thesis reports the former. Eight missing remittance observations were left missing rather than interpolated, seven of them in the group under study. And an earlier version of the projection-sensitivity check produced a convergence result that turned out to be an artefact of a shrinking country sample; rebuilding it on balanced samples reversed the finding.

\subsection{What follows}
\label{sec:further-research}

The clearest implication is about units of analysis. This thesis shows that the Central Asian difference lives in persistent between-country variation and that a country-level design cannot decompose it. The literature that can is micro-level, and it is fragmented: separate surveys, separate countries, separate periods. A harmonised individual-level dataset covering the Central Asian states and a comparison group elsewhere in the former Soviet Union, with ethnicity, parity, marital history and religiosity measured consistently, would allow the mechanisms discussed in Chapter~\ref{sec:discussion} to be tested against one another rather than assembled from separate studies.

Three narrower directions follow from specific limitations. Because period fertility rates confound the number of births with their timing, tempo-adjusted or cohort measures would establish how much of the recent divergence is a change in family size and how much a change in scheduling --- a question this thesis raises repeatedly and cannot settle. Because ethnic composition may account for the majority of the Central Asian fertility increase, and this design has no ethnic dimension at all, that omission is the most consequential one to repair. And the generational reading offered in Chapter~\ref{sec:discussion} --- that the Soviet-era dampening of the pathway from religiosity to fertility weakens for cohorts formed after 1991, and that these are the cohorts now at peak childbearing --- is a conjecture that can be tested directly by interacting religiosity with birth cohort in existing survey data. It is the only mechanism I found that predicts the timing of the divergence rather than just its level, which makes it worth testing.

\subsection{Closing}
\label{sec:closing}

The countries in this study spent most of the twentieth century under one government, and much of the literature on their demography has been organised around the expectation that they would grow more alike once it ended. On fertility they have not. Two of them entered this century with rates 1.4 children apart and left it 2.3 apart, and the widening is recent enough that most of the research on the region predates it.

This thesis measures that divergence, shows that the standard economic explanations do not account for the greater part of it, and marks out the boundary of what fourteen countries observed over twenty-four years can establish. What lies beyond the boundary --- marriage, ethnicity, gender norms, expectations about the future --- is where the answer is, and it is not reachable from here.

%  No tables or figures. No new citations beyond those already used.
% =====================================================================

\section{Conclusion}
\label{sec:conclusion}

\subsection{The two answers}
\label{sec:two-answers}

Fertility in the post-Soviet space did not converge between 2000 and 2023. The four Central Asian countries remained above the other ten in every year of the period without a single crossing, initially higher-fertility countries did not decline faster, and the difference between the two groups ended wider than it began. Beneath a pooled dispersion series that appears flat, Central Asia grew more homogeneous while the rest of the region pulled apart. Around 2016 the pattern changed: dispersion reached its minimum and reversed, and the difference between the groups moved from 1.19 children per woman to 1.65 in seven years. Roughly seventy per cent of that widening came from the comparison countries falling rather than from Central Asia climbing.

The second answer is bounded, in the way set out in Chapter~\ref{sec:introduction}. The Central Asian difference is $+1.314$ children per woman before controls and $+0.964$ after them. Income, urbanisation, remittances and child survival account for about a quarter of it. The remaining three quarters is not explained by this analysis, and the analysis is clear about why: within countries these variables show no stable association with fertility, and across countries the markers that do correlate with the difference --- Muslim population share, age at marriage, female schooling --- cannot be separated from Central Asian regional identity in a sample of fourteen. What the thesis establishes is the size and durability of what remains, not its composition.

\subsection{What is new here}
\label{sec:contribution}

Three things.

The first is coverage. No previous study places all fourteen comparable post-Soviet states in a single frame with consistent measurement, a common window and inference appropriate to the sample size. The comparative post-communist literature largely excludes Central Asia from its formal analysis; the Central Asian literature compares within the region. This thesis puts both halves of the former Soviet Union side by side, and the fact that the difference never closes in any year of the panel is only visible from that vantage point.

The second is timing. The reversal around 2016 and the widening that follows it postdate the windows of essentially all the comparative work on the region. So does the Uzbek fertility surge, the largest single movement in the panel, which took that country from 2.57 to 3.50 children per woman in six years while its population of women in peak childbearing ages was shrinking. The finding that the recent divergence comes mainly from the low-fertility side rather than the high-fertility side is, as far as I am aware, not documented elsewhere, and it changes what needs explaining.

The third is the boundary itself. A null result reported carefully is worth more than a positive result reported loosely, and this thesis reports several. No control variable is individually distinguishable from zero. Within-country variation yields nothing stable. The cross-section cannot separate religion from region. Rather than working around these, the analysis treats them as the finding: the standard macroeconomic account of post-communist fertility, applied with four covariates over twenty-four years, does not reach the Central Asian difference. That is informative precisely because the account is the one the comparative literature relies on.

Some of the methodological choices were made in the same spirit. The wild-cluster bootstrap returns a p-value of 0.033 on the primary estimate where the asymptotic clustered p-value is below 0.01, and the thesis reports the former. Eight missing remittance observations were left missing rather than interpolated, seven of them in the group under study. And an earlier version of the projection-sensitivity check produced a convergence result that turned out to be an artefact of a shrinking country sample; rebuilding it on balanced samples reversed the finding.

\subsection{What follows}
\label{sec:further-research}

The clearest implication is about units of analysis. This thesis shows that the Central Asian difference lives in persistent between-country variation and that a country-level design cannot decompose it. The literature that can is micro-level, and it is fragmented: separate surveys, separate countries, separate periods. A harmonised individual-level dataset covering the Central Asian states and a comparison group elsewhere in the former Soviet Union, with ethnicity, parity, marital history and religiosity measured consistently, would allow the mechanisms discussed in Chapter~\ref{sec:discussion} to be tested against one another rather than assembled from separate studies.

Three narrower directions follow from specific limitations. Because period fertility rates confound the number of births with their timing, tempo-adjusted or cohort measures would establish how much of the recent divergence is a change in family size and how much a change in scheduling --- a question this thesis raises repeatedly and cannot settle. Because ethnic composition may account for the majority of the Central Asian fertility increase, and this design has no ethnic dimension at all, that omission is the most consequential one to repair. And the generational reading offered in Chapter~\ref{sec:discussion} --- that the Soviet-era dampening of the pathway from religiosity to fertility weakens for cohorts formed after 1991, and that these are the cohorts now at peak childbearing --- is a conjecture that can be tested directly by interacting religiosity with birth cohort in existing survey data. It is the only mechanism I found that predicts the timing of the divergence rather than just its level, which makes it worth testing.

\subsection{Closing}
\label{sec:closing}

The countries in this study spent most of the twentieth century under one government, and much of the literature on their demography has been organised around the expectation that they would grow more alike once it ended. On fertility they have not. Two of them entered this century with rates 1.4 children apart and left it 2.3 apart, and the widening is recent enough that most of the research on the region predates it.

This thesis measures that divergence, shows that the standard economic explanations do not account for the greater part of it, and marks out the boundary of what fourteen countries observed over twenty-four years can establish. What lies beyond the boundary --- marriage, ethnicity, gender norms, expectations about the future --- is where the answer is, and it is not reachable from here.
